% Source-numbering note: arXiv:1606.00608 uses one theorem counter for
% definitions, propositions, theorems, examples, and observations in Section 4.
% The three examples after Theorem 4.9 are numbered 4.10--4.12; hence
% `thm:IV.13` is Theorem 4.14 in the arXiv source.

\begin{definition}[Support algebras and blocking maps]%
    \label{def:mpdo_algebra_structure_data}
    \lean{MPOTensor.AlgebraStructureData}
    \leanok
    \uses{def:mpo_tensor}
    Let $\{\mathcal{A}_n\}_{n \ge 0}$ be a family of $*$-subalgebras of
    $\MN{D}$, one at every blocking size, equipped with bilinear maps
    $m_n : \mathcal{A}_n \times \mathcal{A}_n \to \mathcal{A}_{2n}$ and
    linear inclusion maps $\iota_n : \mathcal{A}_n \to \mathcal{A}_{n+1}$.
    These maps are realized in the ambient bond-space matrix algebra by
    \begin{align}
        m_n(x,y)   & = xy, \notag \\
        \iota_n(x) & = x.
        \notag
    \end{align}
\end{definition}

\begin{definition}[Algebras of blocked adjoint fixed points]%
    \label{def:mpdo_blocked_fixed_point_algebra_tower}
    \lean{MPOTensor.HasBlockedAdjointFixedPointAlgebraTower}
    \leanok
    \uses{def:mpdo_algebra_structure_data, def:mpo_tensor}
    An MPO tensor $M$ has a family of algebras of blocked adjoint fixed points
    if there exists a family of support algebras with blocking and inclusion
    maps as in Definition~\ref{def:mpdo_algebra_structure_data} such that, for
    every positive blocking size $n$,
    \begin{align}
        \mathcal{A}_n & = \Fix(E_n^{\dagger}),
        \notag
    \end{align}
    where $E_n$ is the blocked transfer map. This equality is termed
    compatibility with $M$. This condition concerns adjoint fixed points;
    it is not the algebra condition of
    \cite[Theorem~4.14(ii)]{Cirac2016MPDO_arXiv}, which requires a common
    BNT $\{M_\alpha\}_\alpha$, the structure coefficients
    $c_{\alpha,\beta,\gamma}^{(L)}=\tr(\chi_{\alpha,\beta,\gamma}^{L})$,
    and a vector $m$, with $m_\alpha=\tr(\mu_\alpha)$, that is an idempotent
    for the multiplication induced by $c^{(1)}$.
\end{definition}

\begin{theorem}[Idempotent transfer maps and stationary support algebras]%
    \label{thm:mpdo_blocked_fixed_point_algebra_tower_from_zcl}
    \lean{MPOTensor.hasBlockedAdjointFixedPointAlgebraTower_of_isZCL_of_isTP_of_posDef_fixed}
    \leanok
    \uses{def:mpo_is_zcl, def:mpdo_blocked_fixed_point_algebra_tower}
    Assume the transfer map of the doubled tensor is trace-preserving and
    admits a positive-definite fixed point. If $E_1^2 = E_1$, then there is a
    stationary family of support algebras
    \begin{align}
        \mathcal{A}_n & = \Fix(E_1^{\dagger}),
        \notag
    \end{align}
    with multiplication $m_n(x,y) = xy$ and inclusion $\iota_n(x) = x$.
\end{theorem}
\begin{proof}\leanok
    Under the trace-preserving normalization and the faithful fixed point, Wolf's
    fixed-point-algebra theorem gives the $*$-subalgebra
    $\Fix(E_1^{\dagger}) \subseteq \MN{D}$. Since $E_1$ is
    idempotent, $E_n = E_1^n = E_1$ for every $n \ge 1$, so
    $\Fix(E_n^{\dagger}) = \Fix(E_1^{\dagger})$ at every
    positive blocking size, giving the stationary family
    $\mathcal{A}_n = \Fix(E_1^{\dagger})$.
\end{proof}

For the one-site and two-site vertical canonical forms, let
$\{M_\alpha\}_\alpha$ and $\{A_\alpha\}_\alpha$ denote their respective
BNTs. Their boundary contractions have the same local form:
\begin{center}
    \tenkzkernel
    % Formula: M_\alpha(X)_{ij}
    % = \sum_{a,b}(M_\alpha^{ij})_{ab}X_{ba}
    % = \tr(M_\alpha^{ij}X).
    % Source verdict: Papers/1606.00608/DAVID_Fig10.png, cited at source TeX
    % lines 1974--1980, places X immediately left of the boxed M_\alpha.
    % Ink-to-index and contraction audit: the direct adjacent virtual bond and
    % the single-row periodic return below together contract a and b;
    % the unlabeled upper and lower physical stubs carry the free i and j.
    % Paired signature (formula and network):
    % boundary|virtual-west=0|virtual-east=0|physical-up=1|physical-down=1.
    $M_\alpha(X)=$
    \begin{tenkz}[rows={wire}, cols=2, boundary=periodic]
        \tn[skin=box]{X} &
        \tn[skin=mpo, ports={90:physical, 270:physical}]{M_\alpha}
    \end{tenkz}
    \qquad
    % Formula: A_\alpha(X)_{ij}
    % = \sum_{a,b}(A_\alpha^{ij})_{ab}X_{ba}
    % = \tr(A_\alpha^{ij}X).
    % Source verdict: the second panel of DAVID_Fig10.png has the same geometry,
    % with X immediately left of the boxed A_\alpha.
    % Ink-to-index and contraction audit: the direct adjacent virtual bond and
    % the single-row periodic return below together contract a and b;
    % the unlabeled upper and lower physical stubs carry the free i and j.
    % Paired signature (formula and network):
    % boundary|virtual-west=0|virtual-east=0|physical-up=1|physical-down=1.
    $A_\alpha(X)=$
    \begin{tenkz}[rows={wire}, cols=2, boundary=periodic]
        \tn[skin=box]{X} &
        \tn[skin=mpo, ports={90:physical, 270:physical}]{A_\alpha}
    \end{tenkz}
\end{center}
In each cell, $X$ is inserted on the virtual bond and the virtual indices are
closed by $\tr$; the physical legs remain open. In
\cite[Appendix~C.4, lines~1974--1980]{Cirac2016MPDO_arXiv},
the weighted direct sums of these contractions are fixed by compositions
of the reduced tpCPM. Products of these fixed operators are then used to
prove that the compositions are identities on the corresponding direct sums
of matrix algebras. This argument concerns the vertical canonical forms,
not the fixed-point algebras $\Fix(E_n^\dagger)$ of
Definition~\ref{def:mpdo_blocked_fixed_point_algebra_tower}.

\begin{theorem}[Stabilized blocked adjoint fixed points yield support algebras]%
    \label{thm:mpdo_algebra_from_stabilized_adjoint_fixedpoints}
    \lean{MPOTensor.hasBlockedAdjointFixedPointAlgebraTower_of_adjointFixedPoints_eq_of_isTP_of_posDef_fixed}
    \leanok
    \uses{def:mpdo_blocked_fixed_point_algebra_tower}
    Assume the transfer map of the doubled tensor is trace-preserving and
    admits a positive-definite fixed point. If, for every positive blocking
    size $n$,
    \begin{align}
        \Fix(E_n^{\dagger}) & = \Fix(E_1^{\dagger}),
        \notag
    \end{align}
    then there is a compatible family of support algebras.
\end{theorem}
\begin{proof}\leanok
    The trace-preserving map and its positive-definite fixed point give a
    $*$-subalgebra $\Fix(E_1^\dagger)$. Take
    $\mathcal{A}_n := \Fix(E_1^{\dagger})$ for every $n$, with ambient
    multiplication and inclusion. The stabilization hypothesis
    $\Fix(E_n^{\dagger}) = \Fix(E_1^{\dagger})$ then gives
    $\mathcal{A}_n = \Fix(E_n^{\dagger})$ at every positive blocking
    size $n$, which is compatibility with $M$.
\end{proof}

Stabilization of adjoint fixed points does not establish the converse direction of
\cite[Theorem~4.14]{Cirac2016MPDO_arXiv}. In particular,
\cite[Theorem~4.14(ii)]{Cirac2016MPDO_arXiv} requires positive diagonal
matrices $\chi_{\alpha,\beta,\gamma}$ for the multiplication of the vertical
operators $O_L(M_\alpha)$ and a vector $m_\alpha=\tr(\mu_\alpha)$ that is
an idempotent for the multiplication induced by $c^{(1)}$.
By contrast, choosing a basis in each $\mathcal{A}_n$ gives coordinates for
multiplication and inclusion between different blocking sizes. These
coordinates alone do not identify the vertical BNT or its structure
coefficients.

\begin{definition}[Coordinates in the support algebras]%
    \label{def:mpdo_blocked_coefficients}
    \lean{MPOTensor.AlgebraStructureData.toBlockedCoefficients}
    \leanok
    \uses{def:mpdo_algebra_structure_data}
    For each blocking size $n$, fix a basis
    $\{b_i^{(n)}\}_{i \in I_n}$ of $\mathcal{A}_n$ over a finite index set
    $I_n$. The coordinates of $x \in \mathcal{A}_n$ are the unique scalars
    $(x_i)_{i \in I_n} \in \C^{I_n}$ such that
    \begin{align}
        x & = \sum_{i \in I_n} x_i\, b_i^{(n)}.
        \notag
    \end{align}
\end{definition}

\begin{theorem}[Reconstruction from coordinates]%
    \label{thm:mpdo_reconstruct_blocked_coefficients}
    \lean{MPOTensor.AlgebraStructureData.reconstructFromBlockedCoefficients_apply}
    \leanok
    \uses{def:mpdo_blocked_coefficients}
    Reconstructing a coordinate family $(a_i)_{i\in I_n}$ gives
    $\sum_i a_i b_i^{(n)}$ in $\mathcal{A}_n$.
\end{theorem}
\begin{proof}\leanok
    This is the finite-basis reconstruction formula for the chosen basis of
    $\mathcal{A}_n$.
\end{proof}

\begin{definition}[Multiplication coefficients between blocking sizes]%
    \label{def:mpdo_blocked_structure_coefficients}
    \lean{MPOTensor.AlgebraStructureData.blockedStructureCoefficients}
    \leanok
    \uses{def:mpdo_blocked_coefficients}
    For $i,j \in I_n$, define the multiplication coefficients
    $(c^{(n)}_{i,j,k})_{k \in I_{2n}}$ by
    \begin{align}
        m_n(b_i^{(n)},b_j^{(n)})
         & = \sum_{k \in I_{2n}} c^{(n)}_{i,j,k}\, b_k^{(2n)}.
        \notag
    \end{align}
    These are coordinates in $\mathcal{A}_{2n}$, not the same-length
    structure coefficients of the vertical operators $O_L(M_\alpha)$.
\end{definition}

\begin{theorem}[Reconstruction of multiplication from its coefficients]%
    \label{thm:mpdo_mul_coefficients_reconstruct}
    \lean{MPOTensor.AlgebraStructureData.coe_mul_eq_sum_blockedStructureCoefficients}
    \leanok
    \uses{def:mpdo_blocked_structure_coefficients}
    In the ambient bond-space matrix algebra,
    $b_i^{(n)} b_j^{(n)}
        = \sum_{k \in I_{2n}} c^{(n)}_{i,j,k} b_k^{(2n)}$.
\end{theorem}
\begin{proof}\leanok
    Apply the reconstruction formula of
    Theorem~\ref{thm:mpdo_reconstruct_blocked_coefficients} to the element
    $m_n(b_i^{(n)},b_j^{(n)}) \in \mathcal{A}_{2n}$.
\end{proof}

\begin{definition}[Inclusion coefficients between blocking sizes]%
    \label{def:mpdo_blocked_inclusion_coefficients}
    \lean{MPOTensor.AlgebraStructureData.blockedInclusionCoefficients}
    \leanok
    \uses{def:mpdo_blocked_coefficients}
    For $i \in I_n$, define the inclusion coefficients
    $(d^{(n)}_{ij})_{j \in I_{n+1}}$ by
    \begin{align}
        \iota_n(b_i^{(n)})
         & = \sum_{j \in I_{n+1}} d^{(n)}_{ij}\, b_j^{(n+1)}.
        \notag
    \end{align}
\end{definition}

\begin{theorem}[Reconstruction of inclusion from its coefficients]%
    \label{thm:mpdo_inclusion_coefficients_reconstruct}
    \lean{MPOTensor.AlgebraStructureData.reconstructFromBlockedInclusionCoefficients}
    \leanok
    \uses{def:mpdo_blocked_inclusion_coefficients}
    Reconstructing the inclusion coefficients recovers the image of
    $b_i^{(n)}$ under $\iota_n$:
    $\sum_{j \in I_{n+1}} d^{(n)}_{ij} b_j^{(n+1)}
        = \iota_n(b_i^{(n)})$.
\end{theorem}
\begin{proof}\leanok
    \uses{def:mpdo_blocked_coefficients}
    Coefficient extraction and reconstruction are mutually inverse coordinate
    maps for the chosen basis of $\mathcal{A}_{n+1}$. Applying them successively
    to $\iota_n(b_i^{(n)})$ therefore returns the same element.
\end{proof}

\begin{theorem}[Reconstruction of compatible adjoint fixed points]%
    \label{thm:mpdo_fixedpoint_coefficients_reconstruct}
    \lean{MPOTensor.AlgebraStructureData.reconstructFromBlockedCoefficients_of_fixedPoint}
    \leanok
    \uses{def:mpdo_blocked_fixed_point_algebra_tower, def:mpdo_blocked_coefficients}
    If the support algebras are compatible with an MPO tensor $M$ and
    $E_n^\dagger X = X$ for a positive blocking size $n$, then
    $X \in \mathcal{A}_n$. Writing $(X_i)_{i \in I_n}$ for its coordinates
    gives $X = \sum_{i \in I_n} X_i b_i^{(n)}$.
\end{theorem}
\begin{proof}\leanok
    Compatibility identifies $\mathcal{A}_n$ with the adjoint fixed-point algebra of
    $E_n$, so $X$ belongs to $\mathcal{A}_n$.
    Theorem~\ref{thm:mpdo_reconstruct_blocked_coefficients} then reconstructs it.
\end{proof}

\begin{theorem}[Descent of blocked adjoint fixed points]%
    \label{thm:mpdo_adjoint_fix_descent}
    \lean{MPOTensor.adjoint_transferMap_apply_of_hasBlockedAdjointFixedPointAlgebraTower}
    \leanok
    \uses{def:mpdo_blocked_fixed_point_algebra_tower}
    If an MPO tensor has a compatible family of support algebras, then for
    every positive blocking size $n$ and every $X$,
    $E_n^{\dagger}X=X$ implies $E_1^{\dagger}X=X$.
\end{theorem}
\begin{proof}\leanok
    Compatibility at size $n$ places $X$ in $\mathcal{A}_n$. The inclusion
    $\iota_n$ places the same ambient matrix $X$ in $\mathcal{A}_{n+1}$.
    Compatibility at size $n+1$ therefore gives $E_{n+1}^\dagger X=X$.
    Since $E_{n+1}=E_n\circ E_1$, it follows that
    $X=E_{n+1}^\dagger X=E_1^\dagger E_n^\dagger X=E_1^\dagger X$.
\end{proof}

\begin{lemma}[Reverse inclusion of adjoint fixed points]%
    \label{lem:mpdo_adjoint_fix_reverse}
    \lean{MPOTensor.adjoint_blockedTransferMap_apply_of_adjoint_transferMap_apply}
    \leanok
    For every $n$, $E_1^{\dagger}X=X$ implies $E_n^{\dagger}X=X$.
\end{lemma}
\begin{proof}\leanok
    Induct on $n$, using $E_{n+1} = E_n \circ E_1$ and
    $(A \circ B)^{\dagger} = B^{\dagger} \circ A^{\dagger}$.
    The base case is $E_0 = \id$. This argument does not require a compatible
    family of support algebras.
\end{proof}

\begin{lemma}[Blocked powers of adjoint eigenvectors]%
    \label{lem:mpdo_adjoint_blocked_eigenvector}
    \lean{MPOTensor.adjoint_blockedTransferMap_apply_of_adjoint_transferMap_eigenvector}
    \leanok
    If $E=E_1$ and $E^\dagger X=\lambda X$, then
    $E_n^\dagger X=\lambda^n X$ for every $n$.
\end{lemma}
\begin{proof}\leanok
    The case $n=0$ is $E_0=\id$. For the induction step,
    $E_{n+1}=E_n\circ E_1$, hence
    \begin{align}
        E_{n+1}^{\dagger}X
         & = E_1^{\dagger}E_n^{\dagger}X
        = E_1^{\dagger}(\lambda^n X)
        = \lambda^n E_1^{\dagger}X
        = \lambda^{n+1}X.
        \notag
    \end{align}
\end{proof}

\begin{theorem}[Finite-order adjoint eigenvalues under compatibility]%
    \label{thm:mpdo_algebra_root_of_unity_eigenvalue}
    \lean{MPOTensor.adjoint_transferMap_eigenvalue_eq_one_of_hasBlockedAdjointFixedPointAlgebraTower}
    \leanok
    \uses{def:mpdo_blocked_fixed_point_algebra_tower, thm:mpdo_adjoint_fix_descent, lem:mpdo_adjoint_blocked_eigenvector}
    Assume there is a compatible family of support algebras. If $X\ne 0$,
    $E_1^\dagger X=\lambda X$, and $\lambda^n=1$ for some $n>0$, then
    $\lambda=1$.
\end{theorem}
\begin{proof}\leanok
    By Lemma~\ref{lem:mpdo_adjoint_blocked_eigenvector},
    $E_n^\dagger X=\lambda^n X=X$. Theorem~\ref{thm:mpdo_adjoint_fix_descent} gives
    $E_1^\dagger X=X$. Thus $\lambda X=X$, so $(\lambda-1)X=0$. Since
    $X\ne0$, one obtains $\lambda=1$.
\end{proof}

\begin{theorem}[Equality of compatible adjoint fixed-point algebras]%
    \label{thm:mpdo_adjoint_fix_equality_from_algebra}
    \lean{MPOTensor.adjoint_blockedTransferMap_apply_iff_of_hasBlockedAdjointFixedPointAlgebraTower}
    \leanok
    \uses{def:mpdo_blocked_fixed_point_algebra_tower, thm:mpdo_adjoint_fix_descent, lem:mpdo_adjoint_fix_reverse}
    If an MPO tensor has a compatible family of support algebras, then
    $\Fix(E_n^{\dagger})=\Fix(E_1^{\dagger})$ for every positive blocking
    size $n$.
\end{theorem}
\begin{proof}\leanok
    Combine Theorem~\ref{thm:mpdo_adjoint_fix_descent} with
    Lemma~\ref{lem:mpdo_adjoint_fix_reverse}.
\end{proof}

\begin{theorem}[Compatibility of the stationary family of support algebras]%
    \label{thm:mpdo_stationary_tower_from_algebra}
    \lean{MPOTensor.stationaryOfFaithfulFixedPoint_compatible_of_hasBlockedAdjointFixedPointAlgebraTower}
    \leanok
    \uses{def:mpdo_blocked_fixed_point_algebra_tower, thm:mpdo_adjoint_fix_equality_from_algebra}
    Under the trace-preserving normalization and a positive-definite fixed
    point, if $M$ has a compatible family of support algebras, then the
    stationary family $\mathcal{A}_n = \Fix(E_1^\dagger)$ is
    itself compatible with $M$: for every positive blocking size $n$,
    $\Fix(E_1^\dagger)=\Fix(E_n^\dagger)$.
\end{theorem}

\begin{theorem}[Compatible support algebras do not imply an idempotent transfer map]%
    \label{thm:mpdo_algebra_tower_not_imply_zcl}
    \lean{MPOTensor.exists_hasBlockedAdjointFixedPointAlgebraTower_not_isZCL}
    \leanok
    \uses{def:mpdo_blocked_fixed_point_algebra_tower, def:mpo_is_zcl, thm:mpdo_algebra_from_stabilized_adjoint_fixedpoints}
    There exists an MPO tensor of physical and bond dimension two whose
    doubled-tensor transfer map is trace-preserving and has a
    positive-definite fixed point, and which has a compatible family of
    support algebras, although its doubled-tensor transfer map is not
    idempotent.
\end{theorem}
\begin{proof}\leanok
    Take the phase-flip channel with Kraus operators
    $K_0=\frac35 \Id$ and $K_1=\frac45 Z$, where $Z=\diag(1,-1)$.
    Its transfer map satisfies
    \begin{align}
        E(X)_{ij}
         & =
        \begin{cases}
            X_{ij},            & i=j,     \\
            -\frac7{25}X_{ij}, & i\neq j.
        \end{cases}
        \notag
    \end{align}
    The identity is a positive-definite fixed point, and
    $K_0^\dagger K_0+K_1^\dagger K_1=\Id$ proves trace preservation.
    For every $n>0$ and $i\neq j$,
    $E^n(X)_{ij}=(-7/25)^n X_{ij}$. Since $(-7/25)^n\neq 1$ for every $n>0$
    and $E$ is self-adjoint,
    $\Fix((E^n)^\dagger)=\Fix(E^\dagger)$. Thus a compatible family of support
    algebras exists. On the other hand, for $X_{01}\ne0$,
    $E^2(X)_{01}=\frac{49}{625}X_{01}\neq-\frac7{25}X_{01}=E(X)_{01}$.
    Therefore $E^2\neq E$, so the doubled-tensor ZCL condition fails.
\end{proof}

\begin{remark}[Stabilized fixed points and the vertical algebra condition]%
    \label{rem:mpdo_algebra_tower_to_zcl_gap}
    \uses{def:mpdo_blocked_fixed_point_algebra_tower, def:mpo_is_zcl, thm:mpdo_algebra_tower_not_imply_zcl}
    Write $E=E_1$ for the one-site transfer map. Compatibility gives
    $\Fix((E^n)^\dagger)=\Fix(E^\dagger)$ for every $n>0$.
    Hence a nonzero matrix $X$ with $E^\dagger X=\lambda X$ and
    $\lambda^n=1$ lies in $\Fix(E^\dagger)$, so $(\lambda-1)X=0$ and
    $\lambda=1$ by
    Theorem~\ref{thm:mpdo_algebra_root_of_unity_eigenvalue}. Compatibility does
    not imply $E^2=E$, as
    Theorem~\ref{thm:mpdo_algebra_tower_not_imply_zcl} shows. More
    generally, on $\MN{D}$ with $D\ge2$, let $\Pi_{\mathrm{diag}}$ be the
    projection onto diagonal matrices. For $0<\varepsilon<1$, set
    $E(X)=\varepsilon X+(1-\varepsilon)\Pi_{\mathrm{diag}}(X)$, so that
    \begin{align}
        E^n(X)
         & = \Pi_{\mathrm{diag}}(X)
        + \varepsilon^n(X-\Pi_{\mathrm{diag}}(X)),
        \notag                           \\
        E^2-E
         & = (\varepsilon^2-\varepsilon)
        (\id-\Pi_{\mathrm{diag}}).
        \notag
    \end{align}
    The fixed-point algebra is diagonal for every $n>0$, but $E^2\ne E$.

    The coefficient argument in
    \cite[Appendix~C.4]{Cirac2016MPDO_arXiv} instead uses the vertical
    canonical form of an MPDO tensor $M$ in CF. Write its positive diagonal
    multiplicity matrices as $\mu_\alpha$ and its BNT as
    $\{M_\alpha\}_\alpha$. The operators $O_L(M_\alpha)$ are obtained by
    contracting $L$ tensors in the vertical direction, with the original
    physical and virtual roles exchanged. The required same-length product
    law is
    \begin{align}
        O_L(M_\alpha)O_L(M_\beta)
         & =
        \sum_\gamma
        \tr(\chi_{\alpha,\beta,\gamma}^{\,L})
        O_L(M_\gamma),
        \label{eq:mpdo_chi_product_law}
    \end{align}
    where the $\chi_{\alpha,\beta,\gamma}$ are positive diagonal matrices.
    With $m_\alpha=\tr(\mu_\alpha)$, the vector $m$ must be an idempotent
    for the multiplication induced by $c^{(1)}$, namely
    \begin{align}
        m_\gamma
         & =
        \sum_{\alpha,\beta}
        \tr(\chi_{\alpha,\beta,\gamma})
        m_\alpha m_\beta.
        \label{eq:mpdo_idempotent_trace_identity}
    \end{align}
    The proof of \cite[Theorem~4.14]{Cirac2016MPDO_arXiv} uses
    (\ref{eq:mpdo_chi_product_law}) and (\ref{eq:mpdo_idempotent_trace_identity}),
    positivity, and the simple-matrix lemma
    \cite[lines~1155--1163]{Cirac2016MPDO_arXiv} to identify the diagonal
    entries of the two-site multiplicity matrix $\nu_\gamma$ as a multiset:
    \begin{align}
        \{\nu_{\gamma,k}\}_k
         & =
        \{\mu_{\alpha,k_1}\mu_{\beta,k_2}
        \chi_{\alpha,\beta,\gamma,k_3}\}_{\alpha,\beta,k_1,k_2,k_3},
        \notag        \\
        \tr(\nu_\gamma)
         & =m_\gamma.
        \label{eq:mpdo_nu_chi_witness}
    \end{align}
    The reconstruction maps in \cite[lines~2065--2085]{Cirac2016MPDO_arXiv}
    then use
    \begin{align}
        \widetilde R_2\left(\bigoplus_\gamma X_\gamma\right)
         & = \bigoplus_\gamma
        \frac{\nu_\gamma}{m_\gamma}\otimes X_\gamma,
        \notag                \\
        \widetilde R_1\left(\bigoplus_\gamma X_\gamma\right)
         & = \bigoplus_\gamma
        \frac{\mu_\gamma}{m_\gamma}\otimes X_\gamma.
        \label{eq:mpdo_reconstruction_maps}
    \end{align}
    Thus $\tr(\nu_\gamma)=m_\gamma$ and $m_\gamma=\tr(\mu_\gamma)$
    normalize the factors $\nu_\gamma/m_\gamma$ and
    $\mu_\gamma/m_\gamma$, respectively, in
    $T=\widetilde R_2\widetilde T R_1$ and
    $S=\widetilde R_1\widetilde S R_2$. Here $R_1$ and $R_2$ compress to
    the canonical sectors and trace out their multiplicity factors, while
    $\widetilde T$ and $\widetilde S$ conjugate by the block unitaries
    relating the BNTs. These compositions are tpCPM on the retained spaces.
    When discarded zero-sector complements are present, extend the
    compressions by trace-and-prepare branches on the discarded input
    supports to obtain tpCPM on the full physical spaces. These additional
    branches vanish on the tensor letters, so the extended maps still satisfy
    $T[M_1(X)]=M_2(X)$ and $S[M_2(X)]=M_1(X)$, as required for an RFP.

    The scalar Newton--Girard identity relates traces of matrix powers to
    characteristic-polynomial coefficients; see
    the corresponding theorem in the companion quantum-channel volume~\cite[``Newton--Girard trace recursion'']{QICLean2026Blueprint}.
    Such scalar identities do not supply the vertical canonical
    decomposition. The trace-power conclusion of
    \cite[Theorem~4.14(ii)]{Cirac2016MPDO_arXiv} concerns a common BNT and
    positive diagonal matrices satisfying
    (\ref{eq:mpdo_chi_product_law}) and
    (\ref{eq:mpdo_idempotent_trace_identity}). Relating these objects to the
    one-site and two-site vertical canonical forms is required for the maps
    in (\ref{eq:mpdo_reconstruction_maps}); it does not follow from a choice
    of bases in the support algebras. In particular,
    Theorem~\ref{thm:mpdo_algebra_tower_not_imply_zcl} rules out deriving
    idempotence of the doubled-tensor transfer map from compatibility alone.
    % Source-faithful completion requires the BNT-label hypotheses as a
    % separate condition; the support-algebra condition does not imply them.
\end{remark}

\subsection{Diagonal \texorpdfstring{$\chi$}{chi}-matrices and the trace-power formula}

In \cite[Theorem~4.14(ii)]{Cirac2016MPDO_arXiv}, the structure coefficients
of the length-$L$ operator algebra are traces of powers of positive diagonal
matrices $\chi_{\alpha,\beta,\gamma}$. In
Definition~\ref{def:mpdo_diagonal_chi_family}, arbitrary complex diagonal
entries are permitted to separate the scalar trace-power identity from
positivity and operator multiplication.

\begin{definition}[Diagonal \texorpdfstring{$\chi$}{chi} matrices]%
    \label{def:mpdo_diagonal_chi_family}
    \lean{MPOTensor.DiagonalChiFamily}
    \leanok
    Let $I$ be an index set. A family of diagonal complex matrices indexed
    by $(\alpha,\beta,\gamma)\in I^3$ consists of a size
    $r_{\alpha,\beta,\gamma}\in\N$ and entries
    $\chi_{\alpha,\beta,\gamma,k}\in\C$ for
    $0\le k<r_{\alpha,\beta,\gamma}$, defining
    \begin{align}
        \chi_{\alpha,\beta,\gamma}
         & =
        \diag\bigl(    (\chi_{\alpha,\beta,\gamma,k})_{0\le k<r_{\alpha,\beta,\gamma}}
        \bigr).
        \notag
    \end{align}
    A size of zero gives the empty diagonal matrix.
    The matrices $\chi_{\alpha,\beta,\gamma}$ in
    \cite[Theorem~4.14(ii)]{Cirac2016MPDO_arXiv} additionally have positive
    diagonal entries.
\end{definition}

\begin{theorem}[Trace of a power of a diagonal matrix]%
    \label{thm:mpdo_trace_matrix_pow}
    \lean{MPOTensor.DiagonalChiFamily.trace_matrix_pow}
    \leanok
    \uses{def:mpdo_diagonal_chi_family}
    For every triple $(\alpha, \beta, \gamma)$ and every $L \ge 0$,
    \begin{align}
        \tr(\chi_{\alpha,\beta,\gamma}^{L})
         & =
        \sum_{0\le k<r_{\alpha,\beta,\gamma}}
        \chi_{\alpha,\beta,\gamma,k}^{L},
        \notag
    \end{align}
    where $r_{\alpha,\beta,\gamma}$ is the size of
    $\chi_{\alpha,\beta,\gamma}$ and $\chi_{\alpha,\beta,\gamma,k}$ are its
    diagonal entries.
\end{theorem}
\begin{proof}\leanok
    Diagonal matrices raise to powers entrywise and the trace of a diagonal matrix
    is the sum of its entries.
\end{proof}

\begin{definition}[Trace-power form of structure coefficients]%
    \label{def:mpdo_has_chi_trace_power_form}
    \lean{MPOTensor.HasChiTracePowerForm}
    \leanok
    \uses{def:mpdo_diagonal_chi_family}
    Let $c^{(L)}_{\alpha,\beta,\gamma}$ be a family of complex coefficients
    and let $\chi_{\alpha,\beta,\gamma}$ be diagonal matrices on the same
    index set. The coefficients have trace-power form with respect to
    $\chi$ if, for every $L\ge0$ and every triple $(\alpha,\beta,\gamma)$,
    \begin{align}
        c^{(L)}_{\alpha,\beta,\gamma}
         & =
        \tr(\chi_{\alpha,\beta,\gamma}^{L})
        =
        \sum_{0\le k<r_{\alpha,\beta,\gamma}}
        \chi_{\alpha,\beta,\gamma,k}^{L}.
        \notag
    \end{align}
    Here $r_{\alpha,\beta,\gamma}$ is the size of
    $\chi_{\alpha,\beta,\gamma}$. This is the scalar identity appearing in
    \cite[Theorem~4.14(ii)]{Cirac2016MPDO_arXiv}. By itself, the definition
    neither requires positive diagonal entries nor asserts that the
    coefficients describe multiplication of vertical operators.
\end{definition}
